
%%%%%%%%%%%%%%%%%%%%%%%%%%%%%%%%%%%%%%%%%%%%%%%%%%%%%%%%%%%%%%%%%%%%
%                                                                  %
% MINIATURAS MATEM\'ATICAS DE LA GACETA DE LA RSME                 %
%                                                                  %
% $\zeta(2k)$ versus $\zeta(k)$: una relaci�n geom�trica en una espiral  %
% by                                                               %
% Juan Luis Varona                                                 %
%                                                                  %
%%%%%%%%%%%%%%%%%%%%%%%%%%%%%%%%%%%%%%%%%%%%%%%%%%%%%%%%%%%%%%%%%%%%

\documentclass[autocontact]{gaceta}

%---------------------
% La Gaceta se har\'a con esto activado, 
% pero no hace falta ahora si no quieres:
\usepackage[T1]{fontenc}
\usepackage{lmodern}

% Puedes usar esto (u otra codificaci\'on) si quieres:
%\usepackage[latin1]{inputenc}
\usepackage[spanish]{babel}
%\decimalpoint
%---------------------

%---------------------
% Esto ya lo carga el estilo de La Gaceta:
%\usepackage{amsmath, amsthm, amssymb}
%\usepackage{url} % <-- Para p\'aginas web o similar: \url{...}
%\usepackage{graphicx}
%---------------------

%---------------------
% Carga esto para gr\'aficos si lo necesitas:
% \usepackage{wrapfig}
%---------------------

%---------------------
\usepackage{tikz}
\usetikzlibrary{calc}
%---------------------

%---------------------
% <<< ESTO SE AJUSTAR\'A AL EDITAR CADA N\'UMERO DE LA GACETA >>>
%\provisional
\newcounter{pageforminiatura}
\setcounter{pageforminiatura}{352} % Siempre ser\'a p\'agina par
\setcounter{page}{1} % No se usa (dejarlo como 1)
\journame{La Gaceta de la RSME}
\yearofpublication{2015}
\volume{18}
\issuenumber{2}
%---------------------

%---------------------
% Deja esto as\'{\i}:
%\belongstopart{Art\'{\i}culos} 
%\belongstopart{Secciones} 
%\belongstopart{Noticias de la Sociedad} 
%\belongstopart{Actualidad} 
\nopart
%---------------------


% Esto controla el n\'umero de l\'ineas que salen desde el centro
\newcounter{lineas}
\setcounter{lineas}{23}

%---------------------
% Deja esto as\'{\i}: 
\supertitle{\large Miniaturas matem\'aticas de \thejourname\\[1mm]
\hfill\footnotesize
Vol.~\thevolume\ (\theyearofpublication), 
N\'um.~\theissuenumber, P\'ag.~\thepageforminiatura\theprovisional}
%
\noeditor
%---------------------

%---------------------
\title{$\zeta(2k)$ versus $\zeta(k)$: una relaci�n geom�trica en una espiral}
%\author{} % Poner \author{} o no poner nada, da igual.
\contact{Juan Luis Varona, 
Departamento de Matem\'aticas y Computaci\'on, 
Universidad de La Rioja}{jvarona@unirioja.es}
{http://www.unirioja.es/cu/jvarona/}
%---------------------

%%%%%%%%%%%%%%%%%%%%%%%%%%%%%%%%%%%%%%%%%%%%%%%%%%%%%%%%%%%%%%%%%%%%%%%
\begin{document}
%%%%%%%%%%%%%%%%%%%%%%%%%%%%%%%%%%%%%%%%%%%%%%%%%%%%%%%%%%%%%%%%%%%%%%%

%%%%%%%%%%%%%%%%%%%%%%%%%%%%%%%%%%%%%%%%%%%%%%%%%%%%%%%%%%%%%%%%%%%%%%%
% Deja esto as\'{\i}: 
\previousvspace{-1.8cm}
\makesupertitle
\previousvspace{-0.25cm}
\maketitle
\vspace*{-0.4cm}
\thispagestyle{empty} % <<<<< Importante
%%%%%%%%%%%%%%%%%%%%%%%%%%%%%%%%%%%%%%%%%%%%%%%%%%%%%%%%%%%%%%%%%%%%%%%


%\begin{flushleft}
Funci�n zeta de Riemann:
\quad
$\displaystyle \zeta(s) = \sum_{n=1}^{\infty} \frac{1}{n^s},
\quad s \in [1,\infty)$ %% s \ge 1 no refleja que pueda no ser entero
%\end{flushleft}


\smallskip



%%%%%%%%%%%%%%%%%%%%%%%%%%%%%%%%%%%%%%%%%%%%%%%%%%%%%%%%%%% 
\begin{center}
\begin{tikzpicture}[scale=3.6] 
%\clip (-1.3,-0.6666666) rectangle (1.3,1.3); % <-- PARA QUE QUEPA EN UNA P\'AGINA
\coordinate (origin) at (0,0);
\draw[very thin, gray!80] (origin) circle (1);

\coordinate (h1) at (1,0);
% Uso \path en vez de \draw para que no pinte, y pintar al final la espiral
% de una sola vez, pues las esquinas quedan mejor
\path[ultra thick, black] (origin) -- (h1) node[black, midway, below] {$c_1=h_1=1$}; % midway es lo mismo que pos=0.5
\draw (h1) node[above=6pt, left=-2pt] {{\scriptsize $90^\circ$}}; % poner \footnotesize size a 10pt

\coordinate (h2) at ($ (h1) ! 0.5 ! -90:(origin) $); % p. 58
\draw (origin) -- (h2) node[pos=0.5, inner sep=1.5pt, below right] {$h_2$}; % midway es lo mismo que pos=0.5 (p. 193)
\draw (h2) node[above=2.5pt, left=0.5pt] {{\scriptsize $90^\circ$}}; % poner \footnotesize size a 10pt

% \draw (h1) -- (h2) node[midway, right] {$c_2=\dfrac{1}{2^k}$}; % esto pinta la l\'inea
\path (h1) -- (h2) node[midway, right] {$c_2=\dfrac{1}{2^k}$}; % esto no la pinta
% \node [inner sep=1pt, label=right:${c_2=\dfrac{1}{2^k}}$] (c2) at ($ (h1)!.5!(h2) $) {}; % esto no pinta pero define c2

% NOTA: SIN "cm" ES RELATIVO A LA DISTANCIA (h2)--(h3)
\coordinate (h3) at ($ (h2) ! 1cm/3 ! -90:(origin) $); 
\draw (origin) -- (h3) node[pos=0.65, inner sep=1.5pt, below right] {$h_3$}; 
\draw (h3) node[above=0pt, left=0.5pt] {{\scriptsize $90^\circ$}}; % poner \footnotesize size a 10pt

\path (h2) -- (h3) node[midway, right] {$c_3=\dfrac{1}{3^k}$}; % esto no pinta la l\'inea
% \node [inner sep=1pt, label=right:${c_3=\dfrac{1}{3^k}}$] (c3) at ($ (h2)!.5!(h3) $) {}; % esto no pinta pero define c3

\coordinate (h4) at ($ (h3) ! 0.25cm ! -90:(origin) $); 
\draw (origin) -- (h4); 

\newcounter{nmenosuno}
\foreach \n in {5,...,\value{lineas}}
{ %\def\nmenosuno{\n-1} % <------ Esta es la idea
  \pgfmathparse{\n-1}
  \pgfmathsetcounter{nmenosuno}{\pgfmathresult} 
  \coordinate (h\n) at ($ (h\thenmenosuno) ! 1cm/\n ! -90:(origin) $); 
  \draw (origin) -- (h\n);
}

\draw[ultra thick, black] (origin) -- (h1) -- (h2) -- (h3) -- (h4) -- (h5) 
\foreach \n in {5,...,\value{lineas}}
{ 
  -- (h\n)
};

% ...

\newcounter{lineasmasuno}
\newcounter{lineasmastres}

\setcounter{lineasmasuno}{\value{lineas}}
\addtocounter{lineasmasuno}{1} %% o \stepcounter{lineasmasuno} 

\setcounter{lineasmastres}{\value{lineas}}
\addtocounter{lineasmastres}{3}

\foreach \n in {\thelineasmasuno,...,\value{lineasmastres}}
{ 
  \pgfmathparse{\n-1}
  \pgfmathsetcounter{nmenosuno}{\pgfmathresult} 
  \coordinate (h\n) at ($ (h\thenmenosuno) ! 1cm/\n ! -90:(origin) $); 
  \fill[black] (h\n) circle (0.28pt); % (0.25) si la espiral es very thick, y (0.28) si es ultra thick
  \fill ($(origin) ! 0.88 ! (h\n)$) circle (0.2pt);
}

\draw[thin, dotted] (origin) circle ({pi/sqrt(6)});
\end{tikzpicture}
\end{center} 
%%%%%%%%%%%%%%%%%%%%%%%%%%%%%%%%%%%%%%%%%%%%%%%%%%%%%%%%%%% 

\vspace*{-4mm}

\[
\left\{
\begin{aligned}
  &c_n = \dfrac{1}{n^k}, \quad n \ge 1\\
  &h_1 = 1, \quad h_{n}^2 = h_{n-1}^2  + c_{n}^2, \quad n \ge 2
\end{aligned}
\right.
\quad\Rightarrow\quad
h_n = \sqrt{1 + \dfrac{1}{2^{2k}} + \dfrac{1}{3^{2k}} + \cdots + \dfrac{1}{n^{2k}}}
\]

%\[
%  \operatorname{Length}(\text{spiral}) = \zeta(k)
%\]
%\[
%  \operatorname{Radius}(\text{asymptotic circle}) = \sqrt{\zeta(2k)}
%\]

\medskip

\begin{large}
\[
\boxed{%
\begin{gathered}
  \operatorname{Longitud}(\text{espiral}) = \zeta(k) \\
  \operatorname{Radio}(\text{circunferencia asint\'otica exterior}) = \sqrt{\zeta(2k)}\, %%% asymptotic circle
\end{gathered}%
}%
\]
\end{large}


El caso $k=1$ es especialmente llamativo, dado que $\zeta(1)=\infty$ y la espiral da infinitas vueltas acerc\'andose a la circunferencia.

\enlargethispage{2mm}


%%%%%%%%%%%%%%%%%%%%%%%%%%%%%%%%%%%%%%%%%%%%%%%%%%%%%%%%%%%%%%%%%%%%%%%
\end{document}
%%%%%%%%%%%%%%%%%%%%%%%%%%%%%%%%%%%%%%%%%%%%%%%%%%%%%%%%%%%%%%%%%%%%%%%

